\scanpage{43}\mspage{5}

On va définir, pour un système d'iso $u_{ij\unsure{k|h}\unsure{l|\ell}}$, un système de ``rapports anharmoniques''

\[
\rho_{ijkl}:\operatorname{Aut}V_{\unsure{0|i}}\qquad(\{i,j,k,l\}=I\}\tag*{\unsure{$(4_0)$|(4)}}
\]

On a

\[
\begin{aligned}
u_{ijkl}&: V_i\xrightarrow{\sim}V_j\\
u_{ijlk}&: V_i\xrightarrow{\sim}V_j
\end{aligned}
\]

\margin{\unsure{$\times$}\quad\unsure{$'$}\quad .\quad ?\quad .}

d'où

\[
 \rho_{ijkl}=u_{ijlk}^{-1}u_{ijkl}\deleted{\unsure{e}},\tag*{(4)}
\]

\margin{24 quantités}

Donc on aura

\[
 \rho_{ijlk}=(\rho_{ijkl})^{-1}\tag*{(5)}
\]

\margin{\deleted{12} relations \unsure{donc|dont|restent} 12 quantités correspondant aux arêtes orientées de $I$}

On aura d'ailleurs

\[
\underset{\substack{\Vert\\u_{jilk}^{-1}u_{jikl}}}{\rho_{jikl}}=
\deleted{\unsure{u_{ijkl}\rho_{ijkl}}}
\left\{\begin{aligned}
 u_{ij\unsure{kl}}(\rho_{ijkl}^{-1})\deleted{\unsure{u_{ijkl}^{-1}}}\\
 u_{ijlk}(\rho_{ijkl}^{-1})
\end{aligned}\right.\tag*{(6)}
\]

\margin{NB On pose $\beta(\alpha)=\beta\alpha\beta^{-1}$}

[ Ce qui donne

\[
\begin{aligned}
\rho_{j\unsure{,}i,l,k}&=u_{ijkl}(\rho_{ijkl})\\
 &=u_{ijlk}(\rho_{ijkl})
\end{aligned}
\]

\margin{Invariance par double permutation}

% archive sheet number: 42
